Many optical microscopes remain mechanically and optically useful for decades, even when they no longer match the digital and automated workflows expected in modern laboratories or teaching environments. Replacing such instruments can be expensive and wasteful. This thesis investigates whether a vintage microscope can instead be extended with motorized movement, digital imaging, calibration, and automated capturing while preserving the existing microscope body.

The result is RetroScope, an open retrofit platform for motorizing and digitizing vintage microscopes. The prototype attaches motors to the microscope's existing focus and XY-stage controls instead of replacing the whole microscope. A parametric 3D-printable adapter model is used so that important dimensions, such as knob positions and mounting geometry, can be adjusted for other microscope bodies. The system combines low-cost embedded hardware, a custom printed circuit board, a Raspberry Pi, a camera input, and a touchscreen user interface.

The software is designed as a modular application that connects hardware control, live video, calibration, image storage, and automated workflows. Image-based calibration links camera pixels, motor steps, objective profiles, scale information, focus settings, and backlash compensation. This allows the system to support practical workflows such as autofocus, focus stacking, tile scanning, calibrated measurement, and remote access via REST API.

The evaluation shows that the prototype can support these workflows under controlled conditions and that calibration and compensation are essential for making low-cost open-loop mechanics usable. At the same time, the results show clear limitations: the system estimates position from motor commands rather than measuring it directly, and mechanical flex, and backlash still bound the achievable accuracy. RetroScope therefore does not replace a precision research microscope. Instead, it demonstrates a resource-conscious way to give vintage microscopes programmable motion, workflows, and digital capture without discarding the original instrument.
